\documentclass[../../main.tex]{subfiles}
\graphicspath{{../../image/}{../../figure/}}
\begin{document}

\section{可容许弱解及其稳定性性质}\label{section:ch06-可容许弱解稳定性}

在\cref{section:ch04-经典解的崩溃}中我们看到，标量守恒律的初值问题可能容许不止一个弱解，因而需要施加可容许性条件。

在\cref{section:ch04-熵可容性条件}中讨论了熵不等式如何用于该目的。

回顾\cref{section:ch03-例子}，对标量守恒律\eqref{align:ch06-6-1-1}，任一光滑函数 $\eta$ 均可作为熵，相应的熵流为

\begin{align}
Q(u)=\int^{u}\eta^{\prime}(\omega)G^{\prime}(\omega)\,\mathrm{d}\omega,
\label{align:ch06-6-2-1}
\end{align}
且熵产生为零。为方便计，可稍许放松光滑性要求，容许仅局部 Lipschitz 连续的熵（以及相应的熵流）。类似地，$G$ 只需局部 Lipschitz 连续。结果表明，为恰当刻画可容许弱解，须对每个凸熵--熵流对施加熵不等式
\begin{align}
\partial_{t}\eta(u(x,t))+\operatorname{div}Q(u(x,t))\leqslant 0.
\label{align:ch06-6-2-2}
\end{align}

\begin{definition}\label{definition:ch06-6-2-1}
$\mathbb{R}^{m}\times[0,\infty)$ 上的有界可测函数 $u$ 称为\eqref{align:ch06-6-1-1}、\eqref{align:ch06-6-1-2}（初值 $u_{0}\in L^{\infty}(\mathbb{R}^{m})$）的\emph{可容许弱解}，如果对每个凸函数 $\eta$（$Q$ 由\eqref{align:ch06-6-2-1}确定）以及 $\mathbb{R}^{m}\times[0,\infty)$ 上所有具紧支集的非负 Lipschitz 连续试验函数 $\psi$，不等式
\begin{align}
\int_{0}^{\infty}\int_{\mathbb{R}^{m}}\Bigl{[}\partial_{t}\psi\,\eta(u)+\sum_{\alpha=1}^{m}\partial_{\alpha}\psi\,Q_{\alpha}(u)\Bigr{]}\,\mathrm{d}x\,\mathrm{d}t+\int_{\mathbb{R}^{m}}\psi(x,0)\eta(u_{0}(x))\,\mathrm{d}x\geqslant 0
\label{align:ch06-6-2-3}
\end{align}
成立。
\end{definition}

以 $\eta(u)=\pm u$、$Q(u)=\pm G(u)$ 代入\eqref{align:ch06-6-2-3}可知，\eqref{align:ch06-6-2-3}蕴含\eqref{align:ch04-弱解0}，即按\cref{definition:ch06-6-2-1}意义的可容许弱解特别地是\cref{section:ch04-弱解}中定义的弱解。

还注意到，若 $u$ 是\eqref{align:ch06-6-1-1}、\eqref{align:ch06-6-1-2}的经典解，则\eqref{align:ch06-6-2-3}自动以等式成立，即所有经典解均可容许。\eqref{align:ch06-6-2-3}的若干动机将在后续章节给出。

为对所有凸 $\eta$ 验证\eqref{align:ch06-6-2-3}，只需对某族凸 $\eta$ 检验即可，只要其成员的非负系数线性组合张成全体凸函数集合。为给出例子，采用如下标准记号：对 $w\in\mathbb{R}$，$w^{+}$ 表示 $\max\{w,0\}$，$\operatorname{sgn}w$ 在 $w$ 为负、正、零时分别取 $-1$、$1$、$0$。注意任一 Lipschitz 连续函数都是分段线性凸函数列的极限：
\begin{align}
c_{0}u+\sum_{i=1}^{k}c_{i}(u-u_{i})^{+},
\label{align:ch06-6-2-4}
\end{align}
其中 $c_{i}>0$，$i=1,\cdots,k$。因而，只需对熵 $\pm u$（熵流 $\pm G$）以及熵--熵流对族
\begin{align}
\eta(u;\bar{u})=(u-\bar{u})^{+},\qquad Q(u;\bar{u})=\operatorname{sgn}(u-\bar{u})^{+}\,[G(u)-G(\bar{u})],
\label{align:ch06-6-2-5}
\end{align}
验证\eqref{align:ch06-6-2-3}即可，其中 $\bar{u}$ 是取值于 $\mathbb{R}$ 的参数。同样可采用著名的 Kruzkov 熵--熵流对族：
\begin{align}
\eta(u;\bar{u})=|u-\bar{u}|,\qquad Q(u;\bar{u})=\operatorname{sgn}(u-\bar{u})\,[G(u)-G(\bar{u})].
\label{align:ch06-6-2-6}
\end{align}

由注记~4.5.3可知，对任一（不必凸的）熵--熵流对，可容许弱解 $u$ 使分布 $\partial_{t}\eta(u)+\operatorname{div}Q(u)$ 成为测度。


后续章节将用多种方法证明的基本存在唯一性定理如下：

\begin{theorem}\label{theorem:ch06-6-2-2}
对每个 $u_{0}\in L^{\infty}(\mathbb{R}^{m})$，存在\eqref{align:ch06-6-1-1}、\eqref{align:ch06-6-1-2}的唯一可容许弱解 $u$，且
\begin{align}
u(\cdot,t)\in C^{0}([0,\infty);L^{1}_{\mathrm{loc}}(\mathbb{R}^{m})).
\label{align:ch06-6-2-7}
\end{align}
\end{theorem}

下述定理建立了标量守恒律可容许弱解最重要的性质，即 L$^{1}$ 稳定性与 L$^{\infty}$ 单调性：

\begin{theorem}\label{theorem:ch06-6-2-3}
设 $u$ 与 $\bar{u}$ 是\eqref{align:ch06-6-1-1}的可容许弱解，相应初值 $u_{0}$、$\bar{u}_{0}$ 取值于紧区间 $[a,b]$。存在只依赖于 $[a,b]$ 的 $s>0$，使得对任意 $t>0$、$r>0$，
\begin{align}
\int_{|x|<r}[u(x,t)-\bar{u}(x,t)]^{+}\,\mathrm{d}x\leqslant\int_{|x|<r+st}[u_{0}(x)-\bar{u}_{0}(x)]^{+}\,\mathrm{d}x,
\label{align:ch06-6-2-8}
\end{align}
\begin{align}
\|u(\cdot,t)-\bar{u}(\cdot,t)\|_{L^{1}(\mathscr{B}_{r})}\leqslant\|u_{0}(\cdot)-\bar{u}_{0}(\cdot)\|_{L^{1}(\mathscr{B}_{r+st})}.
\label{align:ch06-6-2-9}
\end{align}
此外，若
\begin{align}
u_{0}(x)\leqslant\bar{u}_{0}(x),\qquad\text{a.e.\ 在 }\mathbb{R}^{m}\text{ 上},
\label{align:ch06-6-2-10}
\end{align}
则
\begin{align}
u(x,t)\leqslant\bar{u}(x,t),\qquad\text{a.e.\ 在 }\mathbb{R}^{m}\times[0,\infty)\text{ 上}.
\label{align:ch06-6-2-11}
\end{align}
特别地，$u$ 与 $\bar{u}$ 的（本质）值域均含于 $[a,b]$。
\end{theorem}

\begin{proof}
导出\eqref{align:ch06-6-2-8}的标量守恒律的突出特点在于：由\eqref{align:ch06-6-2-5}定义的函数 $\eta(u;\bar{u})$、$Q(u;\bar{u})$，不仅对固定 $\bar{u}$ 关于变量 $u$ 构成熵--熵流对，而且对固定 $u$ 关于变量 $\bar{u}$ 也构成熵--熵流对。

考虑任一非负 Lipschitz 连续函数 $\phi(x,t,\bar{x},\bar{t})$，定义在 $\mathbb{R}^{m}\times[0,\infty)\times\mathbb{R}^{m}\times[0,\infty)$ 上且具紧支集。固定 $(\bar{x},\bar{t})\in\mathbb{R}^{m}\times[0,\infty)$，对熵--熵流对 $\eta(u;\bar{u}(\bar{x},\bar{t}))$、$Q(u;\bar{u}(\bar{x},\bar{t}))$ 及试验函数 $\psi(x,t)=\phi(x,t,\bar{x},\bar{t})$ 写出\eqref{align:ch06-6-2-3}：
\begin{align}
&\int_{0}^{\infty}\int_{\mathbb{R}^{m}}\Bigl{\{}\partial_{t}\phi(x,t,\bar{x},\bar{t})\,\eta(u(x,t);\bar{u}(\bar{x},\bar{t}))+\sum_{\alpha=1}^{m}\partial_{x_{\alpha}}\phi(x,t,\bar{x},\bar{t})\,Q_{\alpha}(u(x,t);\bar{u}(\bar{x},\bar{t}))\Bigr{\}}\,\mathrm{d}x\,\mathrm{d}t\notag\\
&\qquad+\int_{\mathbb{R}^{m}}\phi(x,0,\bar{x},\bar{t})\,\eta(u_{0}(x);\bar{u}(\bar{x},\bar{t}))\,\mathrm{d}x\geqslant 0.
\label{align:ch06-6-2-12}
\end{align}

交换 $u$ 与 $\bar{u}$ 的地位，类似地对任意固定点 $(x,t)\in\mathbb{R}^{m}\times[0,\infty)$ 得到：
\begin{align}
&\int_{0}^{\infty}\int_{\mathbb{R}^{m}}\Bigl{\{}\partial_{\bar{t}}\phi(x,t,\bar{x},\bar{t})\,\eta(u(x,t);\bar{u}(\bar{x},\bar{t}))+\sum_{\alpha=1}^{m}\partial_{\bar{x}_{\alpha}}\phi(x,t,\bar{x},\bar{t})\,Q_{\alpha}(u(x,t);\bar{u}(\bar{x},\bar{t}))\Bigr{\}}\,\mathrm{d}\bar{x}\,\mathrm{d}\bar{t}\notag\\
&\qquad+\int_{\mathbb{R}^{m}}\phi(x,t,\bar{x},0)\,\eta(u(x,t);\bar{u}_{0}(\bar{x}))\,\mathrm{d}\bar{x}\geqslant 0.
\label{align:ch06-6-2-13}
\end{align}

将\eqref{align:ch06-6-2-12}关于 $(\bar{x},\bar{t})$、\eqref{align:ch06-6-2-13}关于 $(x,t)$ 在 $\mathbb{R}^{m}\times[0,\infty)$ 上积分，再将所得不等式相加，得
\begin{align}
&\int_{0}^{\infty}\int_{\mathbb{R}^{m}}\int_{0}^{\infty}\int_{\mathbb{R}^{m}}\Bigl{\{}(\partial_{t}+\partial_{\bar{t}})\phi(x,t,\bar{x},\bar{t})\,\eta(u(x,t);\bar{u}(\bar{x},\bar{t}))\notag\\
&\qquad+\sum_{\alpha=1}^{m}(\partial_{x_{\alpha}}+\partial_{\bar{x}_{\alpha}})\phi(x,t,\bar{x},\bar{t})\,Q_{\alpha}(u(x,t);\bar{u}(\bar{x},\bar{t}))\Bigr{\}}\,\mathrm{d}x\,\mathrm{d}t\,\mathrm{d}\bar{x}\,\mathrm{d}\bar{t}\notag\\
&\qquad+\int_{0}^{\infty}\int_{\mathbb{R}^{m}}\int_{\mathbb{R}^{m}}\phi(x,0,\bar{x},\bar{t})\,\eta(u_{0}(x);\bar{u}(\bar{x},\bar{t}))\,\mathrm{d}x\,\mathrm{d}\bar{x}\,\mathrm{d}\bar{t}\notag\\
&\qquad+\int_{0}^{\infty}\int_{\mathbb{R}^{m}}\int_{\mathbb{R}^{m}}\phi(x,t,\bar{x},0)\,\eta(u(x,t);\bar{u}_{0}(\bar{x}))\,\mathrm{d}x\,\mathrm{d}\bar{x}\,\mathrm{d}\bar{t}\geqslant 0.
\label{align:ch06-6-2-14}
\end{align}

取 $\mathbb{R}$ 上光滑非负、具紧支集且总质量为一的函数 $\rho$：
\begin{align}
\int_{-\infty}^{\infty}\rho(\xi)\,\mathrm{d}\xi=1.
\label{align:ch06-6-2-15}
\end{align}
考虑 $\mathbb{R}^{m}\times[0,\infty)$ 上任一具紧支集的非负 Lipschitz 试验函数 $\psi$。对充分小的正数 $\varepsilon$，以
\begin{align}
\phi(x,t,\bar{x},\bar{t})=\varepsilon^{-(m+1)}\psi\Bigl{(}\frac{x+\bar{x}}{2},\frac{t+\bar{t}}{2}\Bigr{)}\rho\Bigl{(}\frac{t-\bar{t}}{2\varepsilon}\Bigr{)}\prod_{\beta=1}^{m}\rho\Bigl{(}\frac{x_{\beta}-\bar{x}_{\beta}}{2\varepsilon}\Bigr{)}
\label{align:ch06-6-2-16}
\end{align}
代入\eqref{align:ch06-6-2-14}，再令 $\varepsilon\downarrow 0$。注意到
\begin{align}
(\partial_{t}+\partial_{\bar{t}})\phi(x,t,\bar{x},\bar{t})=\varepsilon^{-(m+1)}\partial_{t}\psi\Bigl{(}\frac{x+\bar{x}}{2},\frac{t+\bar{t}}{2}\Bigr{)}\rho\Bigl{(}\frac{t-\bar{t}}{2\varepsilon}\Bigr{)}\prod_{\beta=1}^{m}\rho\Bigl{(}\frac{x_{\beta}-\bar{x}_{\beta}}{2\varepsilon}\Bigr{)},
\label{align:ch06-6-2-17}
\end{align}
\begin{align}
(\partial_{x_{\alpha}}+\partial_{\bar{x}_{\alpha}})\phi(x,t,\bar{x},\bar{t})=\varepsilon^{-(m+1)}\partial_{\alpha}\psi\Bigl{(}\frac{x+\bar{x}}{2},\frac{t+\bar{t}}{2}\Bigr{)}\rho\Bigl{(}\frac{t-\bar{t}}{2\varepsilon}\Bigr{)}\prod_{\beta=1}^{m}\rho\Bigl{(}\frac{x_{\beta}-\bar{x}_{\beta}}{2\varepsilon}\Bigr{)},
\label{align:ch06-6-2-18}
\end{align}
\begin{align}
|\eta(u(x,t);\bar{u}_{0}(\bar{x}))-\eta(u_{0}(x);\bar{u}_{0}(\bar{x}))|\leqslant|u(x,t)-u_{0}(x)|,
\label{align:ch06-6-2-19}
\end{align}
\begin{align}
|\eta(u_{0}(x);\bar{u}(\bar{x},\bar{t}))-\eta(u_{0}(x);\bar{u}_{0}(\bar{x}))|\leqslant|\bar{u}(\bar{x},\bar{t})-\bar{u}_{0}(\bar{x})|,
\label{align:ch06-6-2-20}
\end{align}
回顾\cref{theorem:ch04-熵凸正则性}并利用标准收敛定理，我们得到

\begin{align}
&\int_{0}^{\infty}\int_{\mathbb{R}^{m}}\Bigl{\{}\partial_{t}\psi(x,t)\,\eta(u(x,t);\bar{u}(x,t))+\sum_{\alpha=1}^{m}\partial_{\alpha}\psi(x,t)\,Q_{\alpha}(u(x,t);\bar{u}(x,t))\Bigr{\}}\,\mathrm{d}x\,\mathrm{d}t\notag\\
&\qquad+\int_{\mathbb{R}^{m}}\psi(x,0)\eta(u_{0}(x);\bar{u}_{0}(x))\,\mathrm{d}x\geqslant 0.
\label{align:ch06-6-2-21}
\end{align}

由\eqref{align:ch06-6-2-5}易见存在 $s>0$，使对解值域中的一切 $u$、$\bar{u}$ 有
\begin{align}
|Q(u;\bar{u})|\leqslant s\,\eta(u;\bar{u}).
\label{align:ch06-6-2-22}
\end{align}

固定 $r>0$、$t\geqslant 0$ 及充分小的 $\varepsilon>0$；以 $\psi(x,\tau)=\chi(x,\tau)\omega(\tau)$ 代入\eqref{align:ch06-6-2-21}，其中 $\chi$、$\omega$ 由\eqref{align:ch05-involution-regularity}与\eqref{align:ch05-cone-convex}定义，得

\begin{align}
\frac{1}{\varepsilon}\int_{t}^{t+\varepsilon}\int_{|x|<r}[u(x,\tau)-\bar{u}(x,\tau)]^{+}\,\mathrm{d}x\,\mathrm{d}\tau
&\leqslant\int_{|x|<r+st}[u_{0}(x)-\bar{u}_{0}(x)]^{+}\,\mathrm{d}x\notag\\
&\quad-\frac{1}{\varepsilon}\int_{0}^{t}\int_{r+s(t-\tau)<|x|<r+s(t-\tau)+\varepsilon}\Bigl{[}s\eta(u;\bar{u})+\frac{Q(u;\bar{u})x}{|x|}\Bigr{]}\,\mathrm{d}x\,\mathrm{d}\tau+O(\varepsilon).
\label{align:ch06-6-2-23}
\end{align}
由\eqref{align:ch06-6-2-22}，\eqref{align:ch06-6-2-23}右端第二个积分为非负。于是令 $\varepsilon\downarrow 0$，回顾\cref{theorem:ch04-熵凸正则性}并利用 $[\,\cdot\,]^{+}$ 为凸函数，即得\eqref{align:ch06-6-2-8}。


在\eqref{align:ch06-6-2-8}中交换 $u$ 与 $\bar{u}$ 的地位，得类似不等式，与\eqref{align:ch06-6-2-8}相加即得\eqref{align:ch06-6-2-9}。

由\eqref{align:ch06-6-2-8}，\eqref{align:ch06-6-2-10}显然蕴含\eqref{align:ch06-6-2-11}。特别地，先取 $\bar{u}_{0}(x)\equiv b$、再取 $u_{0}(x)\equiv a$ 应用此单调性，得 $u(x,t)\leqslant b$、$\bar{u}(x,t)\geqslant a$（a.e.）。交换 $u$ 与 $\bar{u}$ 的地位，可知两解的本质值域均含于 $[a,b]$。于是\eqref{align:ch06-6-2-22}中的 $s$ 只依赖于 $[a,b]$。证明完成。
\end{proof}

由\eqref{align:ch06-6-2-9}立即得到如下唯一性与有限依赖结论：

\begin{corollary}\label{corollary:ch06-6-2-4}
\eqref{align:ch06-6-1-1}、\eqref{align:ch06-6-1-2}至多存在一个可容许弱解。
\end{corollary}

\begin{corollary}\label{corollary:ch06-6-2-5}
可容许弱解在任一点 $(\bar{x},\bar{t})$ 的值只依赖于初值在球 $\mathscr{B}_{s\bar{t}}(\bar{x})$ 上的限制。
\end{corollary}

\eqref{align:ch06-6-2-9}的另一个重要推论是：当初值为局部有界变差时，\eqref{align:ch06-6-1-1}的任一可容许弱解本身也是局部有界变差函数：

\begin{theorem}\label{theorem:ch06-6-2-6}
设 $u$ 是\eqref{align:ch06-6-1-1}的可容许弱解，初值 $u_{0}\in BV_{\mathrm{loc}}(\mathbb{R}^{m})$ 取值于区间 $[a,b]$。则 $u\in BV_{\mathrm{loc}}(\mathbb{R}^{m}\times(0,\infty))$。对任一固定 $t>0$，$u(\cdot,t)\in BV_{\mathrm{loc}}(\mathbb{R}^{m})$，且
\begin{align}
TV_{\mathscr{B}_{r}}u(\cdot,t)\leqslant TV_{\mathscr{B}_{r+st}}u_{0}(\cdot),
\label{align:ch06-6-2-24}
\end{align}
对每个 $r>0$ 成立，其中 $s$ 只依赖于 $[a,b]$。
\end{theorem}

\begin{proof}
记 $\{E_{\alpha},\alpha=1,\cdots,m\}$ 为 $\mathbb{R}^{m}$ 的标准正交基。注意对 $\alpha=1,\cdots,m$，由 $\bar{u}(x,t)=u(x+hE_{\alpha},t)$（$h>0$）定义的函数 $\bar{u}$ 是\eqref{align:ch06-6-1-1}以 $\bar{u}_{0}(x)=u_{0}(x+hE_{\alpha})$ 为初值的可容许弱解。因而由\eqref{align:ch06-6-2-9}，对任意 $t\in(0,T)$，
\begin{align}
\int_{|x|<r}|u(x+hE_{\alpha},t)-u(x,t)|\,\mathrm{d}x\leqslant\int_{|x|<r+st}|u_{0}(x+hE_{\alpha})-u_{0}(x)|\,\mathrm{d}x.
\label{align:ch06-6-2-25}
\end{align}
由于 $u_{0}\in BV_{\mathrm{loc}}(\mathbb{R}^{m})$，\cref{theorem:ch01-BV紧性}与\eqref{align:ch01-difference-quotient}给出 $u(\cdot,t)\in BV_{\mathrm{loc}}(\mathbb{R}^{m})$ 且\eqref{align:ch06-6-2-24}成立。


于是 $\partial_{\alpha}u(\cdot,t)$ 是 Radon 测度，在 $\mathbb{R}^{m}$ 中半径为 $r$ 的球上一致有界，且在紧致时间区间上一致。由于 $u$ 有界，由\cref{theorem:ch01-Gauss-Green}知 $\operatorname{div}G(u(\cdot,t))$ 具有相同性质。

特别地，分布 $\partial_{\alpha}u$ 与 $\operatorname{div}G(u)$ 在 $\mathbb{R}^{m}\times(0,\infty)$ 上是局部有限测度。由于\eqref{align:ch06-6-1-1}在分布意义下成立，$\partial_{t}u$ 也是 $\mathbb{R}^{m}\times(0,\infty)$ 上的测度。因此 $u\in BV_{\mathrm{loc}}(\mathbb{R}^{m}\times(0,\infty))$。证明完成。
\end{proof}

\eqref{align:ch06-6-1-1}的平凡常数解不仅在 L$^{1}$ 中、而且在任一 L$^{p}$ 中稳定。由于 $u$ 可重整化，只需对零解建立 L$^{p}$ 稳定性。

\begin{theorem}\label{theorem:ch06-6-2-7}
设 $u$ 是\eqref{align:ch06-6-1-1}、\eqref{align:ch06-6-1-2}的可容许弱解，初值取值于紧区间 $[a,b]$。存在只依赖于 $[a,b]$ 的 $s>0$，使得对任意 $1\leqslant p\leqslant\infty$、$t\geqslant 0$、$r>0$，
\begin{align}
\|u(\cdot,t)\|_{L^{p}(\mathscr{B}_{r})}\leqslant\|u_{0}(\cdot)\|_{L^{p}(\mathscr{B}_{r+st})}.
\label{align:ch06-6-2-26}
\end{align}
\end{theorem}

\begin{proof}
对 $1\leqslant p<\infty$，考虑凸熵 $\eta(u)=|u|^{p}$，熵流 $Q$ 由\eqref{align:ch06-6-2-1}确定。注意存在与 $p$ 无关的 $s>0$，使
\begin{align}
|Q(u)|\leqslant s\,\eta(u),\qquad u\in[a,b].
\label{align:ch06-6-2-27}
\end{align}
固定 $r>0$、$t\geqslant 0$ 及充分小的 $\varepsilon>0$；对上述熵--熵流对及试验函数 $\psi(x,\tau)=\chi(x,\tau)\omega(\tau)$（$\chi$、$\omega$ 由\eqref{align:ch05-involution-regularity}、\eqref{align:ch05-cone-convex}定义）写出\eqref{align:ch06-6-2-3}，得

\begin{align}
\frac{1}{\varepsilon}\int_{t}^{t+\varepsilon}\int_{|x|<r}|u(x,\tau)|^{p}\,\mathrm{d}x\,\mathrm{d}\tau
&\leqslant\int_{|x|<r+st}|u_{0}(x)|^{p}\,\mathrm{d}x\notag\\
&\quad-\frac{1}{\varepsilon}\int_{0}^{t}\int_{r+s(t-\tau)<|x|<r+s(t-\tau)+\varepsilon}\Bigl{[}s\eta(u)+\frac{Q(u)x}{|x|}\Bigr{]}\,\mathrm{d}x\,\mathrm{d}\tau+O(\varepsilon).
\label{align:ch06-6-2-28}
\end{align}
已知 $u$ 的值域含于 $[a,b]$，故由\eqref{align:ch06-6-2-27}，\eqref{align:ch06-6-2-28}右端第二个积分为非负。令 $\varepsilon\downarrow 0$ 并利用 $|u|^{p}$ 为凸函数，即得\eqref{align:ch06-6-2-26}。证明完成。
\end{proof}

后续章节将给出构造\eqref{align:ch06-6-1-1}、\eqref{align:ch06-6-1-2}可容许弱解的各种方法，并给出\cref{theorem:ch06-6-1-1}的另一证明。

\end{document}
